\subsection*{Monoids}
\label{subsec:monoids}

\begin{definitionbox}{Definition (Monoid)}
\begin{definition}[Monoid]
\label{def:algebraic-monoid}
A \textbf{monoid} is a triple $(A,\star,e)$ where $(A,\star)$ is a semigroup
and $e\in A$ is an identity element:
\[
\forall a \in A,\quad a\star e=e\star a=a.
\]
\end{definition}
\end{definitionbox}

\begin{remark*}[Standard quantified statement]
\[
\operatorname{Monoid}(A,\star,e)
\Longleftrightarrow
\operatorname{Semigroup}(A,\star)
\land
\operatorname{IdentityElement}(e,\star,A).
\]
\end{remark*}

\begin{remark*}[Definition predicate reading]
$(A,\star,e)$ is a monoid exactly when $(A,\star)$ is a semigroup and $e$ is a
two-sided identity for $\star$.
\end{remark*}

\begin{remark*}[Negated quantified statement]
\[
\neg\operatorname{Monoid}(A,\star,e)
\Longleftrightarrow
\neg\operatorname{Semigroup}(A,\star)
\lor
\neg\operatorname{IdentityElement}(e,\star,A).
\]
\end{remark*}

\begin{remark*}[Interpretation]
Monoids add a do-nothing element to associative composition.
\end{remark*}

\begin{remark*}[Examples]
The natural numbers under addition with identity $0$, the natural numbers under
multiplication with identity $1$, and strings under concatenation with the
empty string are monoids.
\end{remark*}

\begin{dependencies}
\begin{itemize}
  \item \hyperref[def:algebraic-semigroup]{Semigroup}
  \item \hyperref[def:identity-element]{Identity Element}
\end{itemize}
\end{dependencies}
